\section{失败模式与修复手册：从课程结论到排障 SOP}

\subsection{失败模式 1：只在自博弈池里高分}
症状是离线对战胜率很高，但一接入真实用户或异构 Agent 池就显著退化。根因通常是策略只学到单均衡盆地。

\begin{importantbox}{修复动作}
引入 cross-population league：每轮训练后都与历史版本、人类行为模型、异构外部策略对战，并把退化样本回流训练。
\end{importantbox}

\subsection{失败模式 2：协作链路过长导致收益反转}
第一类失败暴露的是群体分布错配；第二类失败则发生在系统拓扑本身。即使每个子 Agent 单独看来都有效，串联更多讨论者、批评者与验证者也会累积通信开销和错误传播，因此必须用边际收益曲线而不是组件数量判断架构是否值得保留。
多 Agent 层数增加后，质量提升停滞而时延与成本持续上升，典型于 ``反复互审'' 架构。

\begin{knowledgebox}{简化策略}
\begin{itemize}
    \item 先测单次路由 + 单次验证是否已覆盖 80\% 增益；
    \item 超过两轮协商必须给出边际收益证据；
    \item 对长链任务设置强制 early-stop 与单 Agent 回退。
\end{itemize}
\end{knowledgebox}

\subsection{失败模式 3：通信一致但决策不一致}
看上去 Agent 间交流顺畅，但最终动作互相冲突。说明共享语义层不足，存在 ``同词异义''。

\begin{warningbox}{协议层警报}
只要出现 ``文本看似一致但执行冲突''，优先修协议 schema，不要先盲调模型参数。
\end{warningbox}

\subsection{失败模式 4：文化子群体系统性受损}
在谈判或协作任务里，某些用户群体长期收益更低，说明群体建模与加权目标存在偏差。这是训练目标层面的缺陷，不是单次 prompt 能修复的问题。

\begin{table}[H]
\centering
\begin{tabular}{p{3.0cm}p{4.2cm}p{4.2cm}}
\toprule
\textbf{症状} & \textbf{优先排查项} & \textbf{推荐修复} \\
\midrule
跨群体掉点明显 & 训练分布是否覆盖目标人群 & 扩充人群样本并做分群目标重加权 \\
成本暴涨收益小 & 协作轮数与验证器质量 & 减少协商轮次，提升 verifier 准确率 \\
策略易被利用 & exploitability 探测是否缺失 & 增加 best-response 对抗评测 \\
长任务漂移 & 共享状态与记忆压缩策略 & 强化 state schema 与阶段性重规划 \\
\bottomrule
\end{tabular}
\caption{多 Agent 系统常见故障定位表}
\end{table}

\subsection{从排障到治理：最小化运营闭环}
可把线上治理流程固化为 ``监测 $\rightarrow$ 归因 $\rightarrow$ 回流 $\rightarrow$ 回归''：
\begin{enumerate}
    \item 监测：实时记录跨群体性能、冲突率、回退率；
    \item 归因：区分模型能力不足、协议缺陷、分布漂移；
    \item 回流：将失败轨迹标注并进入增量训练池；
    \item 回归：以固定回归集检查修复是否引入新退化。
\end{enumerate}

这个闭环与 Brown 在课程中强调的思想一致：多 Agent 不是一次性设计，而是持续博弈中的动态系统。

\subsection{本章小结}
多数失败不是 ``模型太弱''，而是 ``目标定义、群体建模、协议治理'' 三者有断层。先修系统边界，再追求模型上限。

